% 底本 PDF pp.90--91，原书第二篇第2.5节，原题160--169。
% 记表面张力为 \alpha；普通液滴一界面，肥皂泡两界面。
\thermocapillaryproblem{160}{81}
求毛细压差：(a) 直径 $d=1.5\,\mathrm{\mu m}$
的水银液滴；(b) 直径 $d=3.0\,\mathrm{mm}$
的肥皂泡，肥皂水溶液表面张力
$\alpha=45\,\mathrm{mN/m}$。
\exsolution
球形单界面液滴的 Laplace 压差为
$\Delta p=2\alpha/R=4\alpha/d$。
(a) 取室温水银表面张力
$\alpha_{\rm Hg}\simeq0.48\,\mathrm{N/m}$，
\[
\boxed{\Delta p_{\rm Hg}=\frac{4\alpha_{\rm Hg}}d
=\frac{4(0.48)}{1.5\times10^{-6}}
\simeq1.28\times10^6\,\mathrm{Pa}
\simeq13\,\mathrm{atm}}.
\]
(b) 肥皂泡有内外\emph{两层}气液界面，
所以压差为 $4\alpha/R=8\alpha/d$：
\[
\boxed{\Delta p_{\rm soap}=\frac{8(0.045)}{0.003}
=120\,\mathrm{Pa}\simeq1.2\times10^{-3}\,\mathrm{atm}}.
\]
\excheck 单层液滴与双层肥皂泡的系数相差2倍；
这是本题最容易遗漏的边界计数。

\thermocapillaryproblem{161}{81}
水银容器底部有直径
$d=70\,\mathrm{\mu m}$ 的圆孔。
水银层的最大厚度是多少，才不会从孔中漏出？
\exsolution
水银对孔壁不润湿。
在临界状态下，孔内下凸液面
可提供最大量级毛细压差
$\Delta p_{\max}=4\alpha_{\rm Hg}/d$，
用于抵消液柱静压 $\rho_{\rm Hg}gh$：
\[
\rho_{\rm Hg}gh_{\max}=\frac{4\alpha_{\rm Hg}}d.
\]
因此
\[
\boxed{h_{\max}=\frac{4\alpha_{\rm Hg}}
{\rho_{\rm Hg}gd}}.
\]
取 $\alpha_{\rm Hg}\simeq0.48\,\mathrm{N/m}$、
$\rho_{\rm Hg}\simeq1.36\times10^4\,\mathrm{kg/m^3}$，
\[
\boxed{h_{\max}\simeq0.21\,\mathrm m=21\,\mathrm{cm}}.
\]
\excheck 采用理想的最大曲率阻挡模型；
真实孔缘粗糙程度与接触角滞后
会影响临界漏液厚度。

\thermocapillaryproblem{162}{81}
装有压强为 $p_0$ 空气的容器里有直径为 $d$
的肥皂泡。容器气压在等温条件下减为原来的
$n$ 分之一，原书题面又说肥皂泡直径增大
$n$ 倍。求肥皂水溶液的表面张力。
\exsolution
设泡膜两侧分别为泡内空气和容器空气。
若容器压力从 $p_0$ 降至 $p_0/n$，
泡的直径一般从 $d$ 变为
$\eta d$，两因子\emph{未必相等}。
两膜面的 Laplace 压差为
$8\alpha/d$，故初末泡内气压为
\[
p_i=p_0+\frac{8\alpha}{d},\qquad
p_f=\frac{p_0}{n}+\frac{8\alpha}{\eta d}.
\]
等温气体的量不变，按
$p_id^3=p_f(\eta d)^3$，
得到
\[
p_0+\frac{8\alpha}d
=\frac{p_0\eta^3}n+
\frac{8\alpha\eta^2}d.
\]
因此一般表达式
\[
\boxed{\alpha=
\frac{p_0d}{8}\,
\frac{1-\eta^3/n}{\eta^2-1}}.
\]
\textbf{底本条件冲突：}
若严格把题面中两个倍数都取为 $n>1$
（即 $\eta=n$），则
$\alpha=-p_0d/8<0$，这与普通肥皂水的
正表面张力矛盾。
\excheck 本稿不把第二个 $n$ 悄悄改成另一个
数值：原版书末的通式含独立的 $\eta$，
却与题面把两个倍数都印成 $n$ 的表述不一致；
需查更可靠版本核定第二个倍数，当前保留
\emph{条件性通式}，不虚构正数答案。

\thermocapillaryproblem{163}{81}
水中深度 $h=5.0\,\mathrm m$ 处，
有一个直径 $d=4.0\,\mathrm{\mu m}$
的空气泡。外界大气压为标准值 $p_0$。
求泡内空气的压强。
\exsolution
泡外的水静压为
$p_{\rm out}=p_0+\rho gh$。
空气泡在水中只有一层气液界面，
气泡内压比水侧高
$\Delta p=2\alpha/R=4\alpha/d$。
因此
\[
\boxed{p_{\rm in}=p_0+\rho gh+\frac{4\alpha}{d}}.
\]
取水表面张力
$\alpha\simeq0.073\,\mathrm{N/m}$、
$\rho=1000\,\mathrm{kg/m^3}$，
\[
p_{\rm in}\simeq1.013\times10^5+
(1000)(9.81)(5)+\frac{4(0.073)}{4\times10^{-6}}
\simeq2.23\times10^5\,\mathrm{Pa},
\]
\[
\boxed{p_{\rm in}\simeq2.2\,\mathrm{atm}}.
\]
\excheck 水中气泡是\emph{单界面}，
不能照搬肥皂泡两层界面的 $8\alpha/d$。

\thermocapillaryproblem{164}{81}
池塘底部形成直径 $d=4.0\,\mathrm{\mu m}$
的气泡。当它等温上升至水面时，
直径扩大 $\eta=1.1$ 倍。
标准大气压为 $p_0$，求池塘深度 $h$。
\exsolution
底部水深 $h$ 处泡内压强
$p_b=p_0+\rho gh+4\alpha/d$。
水面处泡直径 $\eta d$，
泡内压强
$p_t=p_0+4\alpha/(\eta d)$。
泡内气体物质的量和温度不变，
满足
$p_bV_b=p_tV_t$，
且 $V_t/V_b=\eta^3$。
故
\[
p_0+\rho gh+\frac{4\alpha}{d}
=\eta^3p_0+\eta^2\frac{4\alpha}{d}.
\]
整理：
\[
\boxed{h=\frac{
p_0(\eta^3-1)+(4\alpha/d)(\eta^2-1)
}{\rho g}}.
\]
代入 $\eta=1.1$、水表面张力约
$0.073\,\mathrm{N/m}$ 得
\[
\boxed{h\simeq5.0\,\mathrm m}.
\]
\excheck 同时保留了水静压和曲面附加压差；
若把气泡看成无表面张力的大气泡，
将得到不同的深度。

\thermocapillaryproblem{165}{81}
两根相互连通的竖直毛细管中装有水银，
内径分别为
$d_1=0.50\,\mathrm{mm}$、
$d_2=1.00\,\mathrm{mm}$，
水银与管壁的接触角为
$\theta=138^\circ$。
求两管水银柱液面的高度差。
\exsolution
圆毛细管的毛细升降量为
$h(d)=4\alpha\cos\theta/(\rho gd)$，
水银不润湿玻璃，$\cos138^\circ<0$，
因此两液面均低于外部参考水面。
高度差的模为
\[
\boxed{\Delta h=
\frac{4\alpha|\cos\theta|}{\rho g}
\left|\frac1{d_1}-\frac1{d_2}\right|
=\frac{4\alpha|\cos\theta|(d_2-d_1)}
{\rho gd_1d_2}}.
\]
取 $\alpha_{\rm Hg}\simeq0.48\,\mathrm{N/m}$、
$\rho_{\rm Hg}\simeq1.36\times10^4\,\mathrm{kg/m^3}$，
得
\[
\boxed{\Delta h\simeq1.1\times10^{-2}\,\mathrm m
=11\,\mathrm{mm}}.
\]
\excheck 水银\emph{下降}而非上升；
题目求的是两管液面差，取绝对值。

\thermocapillaryproblem{166}{81}
一根内径为 $0.50\,\mathrm{mm}$ 的竖直毛细管
插入水中，其高出外水面的部分长
$h=25\,\mathrm{mm}$。
求顶部水面弯月面的曲率半径。
\exsolution
对润湿玻璃的水，上升水柱高度为 $h$。
在顶部弯月面，
Laplace 压差平衡重力压头：
\[
\frac{2\alpha}{R_c}=\rho gh.
\]
因此
\[
\boxed{R_c=\frac{2\alpha}{\rho gh}}.
\]
用 $\alpha\simeq0.073\,\mathrm{N/m}$、
$\rho\simeq1000\,\mathrm{kg/m^3}$、
$h=0.025\,\mathrm m$：
\[
\boxed{R_c\simeq5.95\times10^{-4}\,\mathrm m
=0.60\,\mathrm{mm}}.
\]
\excheck $R_c>d/2$，说明实际弯月面无需
达到理想完全润湿的最大曲率就能支撑
$25\,\mathrm{mm}$ 的水柱。

\thermocapillaryproblem{167}{81}
一根长度 $\ell=110\,\mathrm{mm}$、
内径 $d=20\,\mathrm{\mu m}$ 的玻璃毛细管，
上端密封，下端竖直插入水中。
外界压强为标准大气压。
问毛细管应浸入水中多长 $x$，
才能使管内外水面高度相同？
\exsolution
设密封端最初管内的空气占据长度 $\ell$、
压强为 $p_0$。管内水面最终与外水面齐平时，
进入管内的水柱长度等于浸入深度 $x$，
剩余空气长度为 $\ell-x$。
按等温玻意耳定律，
\[
p_{\rm air}(\ell-x)=p_0\ell,
\quad p_{\rm air}=\frac{p_0\ell}{\ell-x}.
\]
在同一液面高度，管内水侧为 $p_0$；
润湿弯月面使管内空气压强比其高
$4\alpha/d$，所以
\[
\frac{p_0\ell}{\ell-x}-p_0=\frac{4\alpha}{d}.
\]
整理：
\[
\boxed{x=\frac{\ell}{1+p_0d/(4\alpha)}}.
\]
取 $\alpha_{\rm water}=0.073\,\mathrm{N/m}$，
\[
\boxed{x\simeq0.014\,\mathrm m=1.4\,\mathrm{cm}}.
\]
\excheck 气体压缩需满足 $0<x<\ell$；
原书一处排版将“110 mm and”连在一起，
本稿已按图像分词恢复。

\thermocapillaryproblem{168}{82}
一根长度为 $\ell$、上端密封的竖直毛细管
下口刚接触液面，液体上升高度为 $h$。
液体密度为 $\rho$、管内径 $d$、
接触角为 $\theta$，外界压强为 $p_0$。
求液体表面张力 $\alpha$。
\exsolution
设管内空气保持等温且初始充满全长 $\ell$，
液体上升 $h$ 后，空气长度为 $\ell-h$，
所以
$p_{\rm air}=p_0\ell/(\ell-h)$。
管内弯月面下方的液体压强
$p_{\rm liq}=p_0-\rho gh$。
圆管弯月面的 Laplace 压差为
\[
p_{\rm air}-p_{\rm liq}
=\frac{4\alpha\cos\theta}{d}.
\]
将上述两压强代入：
\[
\frac{4\alpha\cos\theta}{d}
=\frac{p_0\ell}{\ell-h}-p_0+\rho gh
=\frac{p_0h}{\ell-h}+\rho gh.
\]
因此
\[
\boxed{\alpha=
\frac d{4\cos\theta}
\left(\rho gh+\frac{p_0h}{\ell-h}\right)}.
\]
\excheck 这一式在 $h\to0$ 时趋于零，
符合无液柱上升的极限。原版书末抽取结果似乎
遗漏了“$-p_0$”一项，需以清晰原图复核；
本稿根据压强平衡逐项列出推导，不直接抄录可疑公式。

\thermocapillaryproblem{169}{82}
一根直径 $d_1=1.5\,\mathrm{mm}$ 的玻璃棒，
同轴插入内径 $d_2=2.0\,\mathrm{mm}$
的玻璃毛细管中。
将装置竖直插入水面，问水在环形间隙中
上升多高？
\exsolution
取完全润湿玻璃 $\theta\simeq0$。
液柱截面积为
$A=\pi(d_2^2-d_1^2)/4$。
两个接触周线的总长度为
$L=\pi(d_1+d_2)$，
表面张力竖直向上的合力
$F_\alpha=\alpha L$。
与上升液柱重力
$\rho gh A$ 平衡：
\[
\alpha\pi(d_1+d_2)
=\rho gh\frac{\pi}{4}(d_2^2-d_1^2).
\]
约去 $d_1+d_2$：
\[
\boxed{h=\frac{4\alpha}{\rho g(d_2-d_1)}}.
\]
代入水的表面张力
$\alpha\simeq0.073\,\mathrm{N/m}$，
\[
\boxed{h\simeq0.060\,\mathrm m=6.0\,\mathrm{cm}}.
\]
\excheck 关键间隙宽度为 $d_2-d_1$；
不能把整个外管直径 $d_2$ 直接代进单管公式。
